\section{Ranging Sensor}

Súčasťou výbavy manipulátora je dištančný snímač (angl. Ranging Sensor), ktorý plní funkciu bezkontaktného merania vzdialenosti k objektom v pracovnom priestore. Na tento účel bol zvolený senzor \textbf{VL53L1X} od spoločnosti STMicroelectronics. 

\subsection{Špecifikácia a vlastnosti senzora VL53L1X}
Senzor VL53L1X funguje na pokročilom princípe Time-of-Flight (ToF). Vyžaruje neviditeľný infračervený laserový lúč a meria absolútny čas letu fotónov (od vyslania po návrat), čím poskytuje presnú vzdialenosť úplne nezávisle od farby či povrchovej odrazivosti prekážky.
Medzi jeho kľúčové vlastnosti patrí:
\begin{itemize}
    \item Maximálny dosah snímania až do 4 metrov v závislosti od podmienok a režimu.
    \item Tri prepínateľné módy (Short, Medium, Long), ktoré dovoľujú optimalizovať senzor na kratšie, no rýchlejšie merania, alebo pomalšie merania do väčších vzdialeností.
    \item Úsporná a jednoduchá sériová komunikácia prostredníctvom I2C zbernice.
\end{itemize}

Snímač je fyzicky upevnený priamo na \textbf{koncovom efektore} robotického ramena. Vďaka tomuto umiestneniu na samom konci kinematického reťazca sa zorné pole senzora posúva a 
natáča spolu s efektorom. To otvára možnosti pre dynamické merania, napríklad 3D mapovanie pracovného priestoru.

\subsection{Softvérová implementácia}
Nižšie je uvedený prehľad implementácie triedy \textbf{RangingSensor}, ktorá slúži ako konfiguračný a potvrdzovací wrapper práve pre modul VL53L1X, pričom rozširuje jeho základnú funkcionalitu o priestorové mapovanie a logovanie dát na SD kartu.

\subsubsection{Čo bolo implementované}

\begin{itemize}
	\item Wrapper trieda pre senzor VL53L1X.
	\item Podpora 3D mapovania priestoru.
	\item Štruktúra SensorPose (x, y, z, roll, pitch, yaw).
	\item Štruktúra MeasurementData (meranie + aktuálna poloha).
	\item SPI inicializácia SD karty.
	\item Zápis nameranych dát do súboru pointcloud.ply pre neskorší rekonštrukciu mapy.
\end{itemize}

\subsection{Verejné funkcie triedy}

\subsubsection{Konštruktor}

\texttt{RangingSensor(TwoWire\& wire)}

Pripraví inštanciu senzora naviazanú na konkrétnu I2C zbernicu (napr. Wire3).

\subsubsection{Inicializácia}

\texttt{bool init(uint32\_t sample\_period\_ms, DistanceMode mode)}

Inicializuje senzor, nastavý časový rozpočet a režim merania.

\begin{itemize}
	\item \texttt{sample\_period\_ms} -- perióda merania v milisekundách.
	\item \texttt{mode} -- režim senzora: Short / Medium / Long.
	\item Návratová hodnota: true pri úspechu, false pri chybe.
\end{itemize}

\subsubsection{Kontinuálne meranie}

\texttt{void startContinuous() / void stopContinuous()}

Spustenie alebo zastavenie kontinuálneho snímania vzdialenosti.

\subsubsection{SD karta a logovanie}

\texttt{bool initSD(uint32\_t cs\_pin)}

Inicializuje SD kartu cez SPI rozhranie.

\texttt{void enableLogging(bool enable, bool clear = false)}

Zapína alebo vypína zápis do súboru pointcloud.ply. Obsahuje buffer manažment (flush\_counter).

\begin{itemize}
	\item \texttt{enable} -- true zapne logovanie, false vypne logovanie.
	\item \texttt{clear} -- true vymaže predchádzajúci súbor.
\end{itemize}

\subsubsection{Získavanie meraní}

\texttt{MeasurementData read(float x, float y, float z, float roll, float pitch, float yaw)}

Načítá najnovšiu vzdialenosť a priradí k nej aktuálnu polohu/orientáciu senzora. Pri zapnutom logovaní automaticky prepočíta lúč do 3D tela a uloží výsledok na SD kartu.

\subsection{Implementácia v robot\_arm\_servo.ino}

\begin{enumerate}
	\item Vytvorenie objektu: \texttt{RangingSensor rangingSensor(Wire3);}
	\item Inicializácia senzora v setup(): 
	
    \texttt{rangingSensor.init(Ts\_ms, RangingSensor::DistanceMode::Short);}
	\item Inicializácia SD karty: \texttt{rangingSensor.initSD(PE4);}
	\item Zapnutie logovania a vymazanie starých dát: 
    
    \texttt{rangingSensor.enableLogging(true, true);}
	\item Spustenie kontinuálneho režimu: \texttt{rangingSensor.startContinuous();}
	\item V hlavnej slučke čítanie meraní: 
    
    \texttt{MeasurementData measurement = rangingSensor.read(x, y, z, roll, pitch, yaw);}
\end{enumerate}

\subsection{Záver}

Trieda RangingSensor zabezpečuje kompletnú integráciu senzora VL53L1X do robotického ramena. Spája meranie vzdialenosti, informácie o polohe a orientácii ramena a vytvárajúci dátový základ pre 3D rekonštrukciu priestoru.
